\documentclass[12pt]{article}

\usepackage[a4paper, margin=1in]{geometry}
\usepackage{amsmath,amssymb,amsthm}
\usepackage{mathtools}
\usepackage{mathrsfs}
\usepackage{enumitem}
\usepackage{xcolor}
\usepackage[hidelinks]{hyperref}
\usepackage{fontspec}

\defaultfontfeatures{Ligatures=TeX}
\IfFontExistsTF{Times New Roman}
  {\setmainfont{Times New Roman}}
  {\setmainfont{TeX Gyre Termes}}
\IfFontExistsTF{Menlo}
  {\setmonofont{Menlo}}
  {\setmonofont{Latin Modern Mono}}

\newtheorem{theorem}{Theorem}[section]
\newtheorem{proposition}[theorem]{Proposition}
\newtheorem{lemma}[theorem]{Lemma}
\newtheorem{corollary}[theorem]{Corollary}
\theoremstyle{definition}
\newtheorem{definition}[theorem]{Definition}
\theoremstyle{remark}
\newtheorem{remark}[theorem]{Remark}

\newcommand{\tightlist}{%
  \setlength{\itemsep}{0pt}\setlength{\parskip}{0pt}}

\title{%
  \textbf{The Structural Horizon of Self-Description:}\\
  \textbf{Fisher Bounds, Small-Ball Remainders, and the Quantitative Explanatory Gap}%
}

\author{
  Sidong Liu, PhD \\
  \small iBioStratix Ltd \\
  \small \texttt{sidongliu@hotmail.com}
}

\date{April 2026}

\begin{document}
\emergencystretch=2em
\raggedbottom
\hbadness=5000
\vbadness=5000

\maketitle

\begin{abstract}
%
Paper~G located the explanatory gap at a structural limit of
self-description, but did so qualitatively. Paper~II then identified a
concrete physical carrier of that limit: the higher-order
\(O(R^{d+3})\) remainder in the small-ball modular expansion, beyond
the part controlled by linearized gravitational consistency. The present
paper supplies the quantitative sequel. We define the self-description
deficit of a finite observer as the optimal local minimax
Araki-relative-entropy loss in compressing and reconstructing its own
accessible state by means of an admissible \(n\)-dimensional self-model.

Our first main result is a local minimax drift-floor theorem. For
locally bounded self-models that factor through a Fisher-whitened
\(n\)-dimensional statistical manifold, the local deficit satisfies
\[
\varepsilon_{\mathrm{loc}}(n)
\ge
\frac{(1-2\eta)\,\mathcal{I}_F}
{2\bigl(\mathcal{I}_F + \mathcal{I}_{\mathrm{ego}}\bigr)},
\]
where \(\mathcal{I}_F\) is the observer's modular Fisher sensitivity,
\(\mathcal{I}_{\mathrm{ego}}\) is the ego-rigidity prior of
T-DOME~Paper~III, and \(\eta\in[0,\tfrac{1}{2})\) is a local
quadratic-control parameter that can be made arbitrarily small by
shrinking the BKM ball. Our second main result then connects the finite
trace-budget side of the theory to Paper~II by defining the maximal
small-ball control order \(k(n)\). In the free scalar CFT worked
example, an exact \(d=1\) signal-factor calculation together with a
general-dimensional boundary-kernel estimate yields
\[
\log I_k = -\alpha k\log k + O(k),
\]
so that
\[
k(n)=O\!\left(\frac{\log n}{\log\log n}\right).
\]
Thus a finite self-model can internalize only finitely many higher-order
small-ball coefficients, and the number of controllable orders grows
only logarithmically.

The paper does not claim a fully operator-algebraic derivation of the
Fisher ceiling, nor a general theorem for arbitrary Type~III nets. Its
claims are deliberately narrower. Within the local regime stated above,
however, the conclusion is sharp: the residual explanatory gap admits a
quantitative reinterpretation as a structural horizon of
self-description.
 \end{abstract}

\clearpage
\tableofcontents
\clearpage

%
\section{Introduction}\label{sec:1}

Paper~G\cite{PaperG} argued that the apparent explanatory gap is not,
at bottom, a failure to explain how physical structure produces
experience. The deeper obstruction is structural. A framework that
describes the world by selecting accessible structure cannot, without
circularity, fully internalize the conditions of its own accessibility.
That result was qualitative. It identified a boundary, but did not
measure it.

Paper~II\cite{PaperII} moved the discussion to a more concrete level.
In the higher-dimensional CHM ball framework, consistency between local
modular data and local geometry fixes the linearized Einstein tensor at
order \(R^{d+1}\). The first uncontrolled terms appear at
\(O(R^{d+3})\). This naturally suggests a new question: can those
higher-order remainders be interpreted as the quantitative residue of
finite self-description?

The present paper gives a first affirmative answer. Its aim is narrower
than a full theory of finite observers, but sharper than a philosophical
gloss. We define the \emph{self-description deficit} of an observer of
complexity \(n\) as a local minimax Araki-relative-entropy loss incurred
when the observer compresses and reconstructs its own accessible state
through an admissible \(n\)-dimensional self-model. We then derive a
local lower bound on this deficit from three inputs:
\begin{enumerate}
\item the BKM/QFIM structure already used in Paper~I\cite{PaperI};
\item the self-referential Fisher and ego-rigidity formalism of
T-DOME~Paper~III\cite{TDOMEIII};
\item the small-ball remainder structure isolated in Paper~II
\cite{PaperII}.
\end{enumerate}

The main theorem chain is intentionally honest. We do \emph{not} claim
that every finite observer automatically determines a finite statistical
model in a purely operator-algebraic way. Nor do we claim a general
Type~III theorem for arbitrary nets or arbitrary perturbations. The
theorem-grade result is conditional on an admissible self-model
factoring through a Fisher-whitened \(n\)-dimensional statistical
manifold. Within that domain, the drift-channel bound for locally
bounded self-models is clean (Corollary~\ref{cor:3.4}, after shrinking
\(\delta\) so that \(\eta\to 0\)):
\[
\varepsilon_{\mathrm{loc}}(n)
\ge
\frac{\mathcal{I}_F}{2\bigl(\mathcal{I}_F+\mathcal{I}_{\mathrm{ego}}\bigr)}.
\]
This already quantifies the structural claim of Paper~G at the local
minimax level.

The second half of the paper translates that bound back into the
small-ball language of Paper~II. We define the maximal small-ball order
\(k(n)\) that a finite Fisher budget can control. In the free scalar
CFT worked example, the signal factor has exact factorial suppression in
\(d=1\), while the general-dimensional renormalized boundary kernel
contributes at most polynomial \(k\)-dependence. The resulting Fisher
spectrum obeys
\[
\log I_k = -\alpha k\log k + O(k),
\]
and therefore
\[
k(n)=O\!\left(\frac{\log n}{\log\log n}\right).
\]
The controllable part of the small-ball expansion grows only
logarithmically with model complexity.

The paper is organized as follows. Section~\ref{sec:2} defines finite
observers, admissible self-models, and the self-description deficit.
Section~\ref{sec:3} proves the local drift-floor theorem. Section~\ref{sec:4}
defines the small-ball control order and derives its abstract growth
law from a spectral hypothesis. Section~\ref{sec:5} verifies that
hypothesis at worked-example level in the free scalar CFT. Section~\ref{sec:6}
explains the trilogy-level significance of the result, and
Section~\ref{sec:7} states the remaining limitations and open problems.

\section{Finite Observers and the Self-Description Deficit}\label{sec:2}

\subsection{Finite Observers}\label{subsec:2.1}

Paper~G\cite{PaperG} framed structural incompleteness in architectural
terms: no accessibility-based framework can fully internalize its own
accessibility conditions without either regress or collapse. To
quantify that boundary, we must first specify what counts as a finite
self-model.

\begin{definition}[Finite observer]\label{def:2.1}
A finite observer of complexity \(n\) is a triple
\[
(\mathcal{A}_{\mathrm{obs}},\omega_{\mathrm{obs}},\mathfrak{M}_n),
\]
where:
\begin{enumerate}
\item \(\mathcal{A}_{\mathrm{obs}} \subset \mathcal{A}_c\) is the
observer's accessible subalgebra;
\item \(\omega_{\mathrm{obs}} = \omega|_{\mathcal{A}_{\mathrm{obs}}}\)
is the restricted state;
\item \(\mathfrak{M}_n\) is an \(n\)-dimensional statistical model used
internally to represent the observer's own accessible state.
\end{enumerate}
\end{definition}

The complexity parameter \(n\) measures effective statistical degrees of
freedom, not arbitrary description length. This is deliberate. Paper~III
is not about Kolmogorov complexity in full generality. It is about the
finite number of statistically usable directions available to an
observer trying to model itself.

\subsection{Modular Fisher Sensitivity}\label{subsec:2.1b}

The Tomita--Takesaki modular automorphism group
\(\sigma_t^\omega\) of a faithful normal state \(\omega\) on a von
Neumann algebra satisfies \(\omega\circ\sigma_t^\omega=\omega\) for
all \(t\). The orbit is therefore trivial at the level of states, and
the associated tangent vector vanishes. Since the modular orbit itself
carries no state-level Fisher information, we replace it with the
intrinsic accessible-unitary perturbation family as the operational
proxy for modular sensitivity.

\begin{definition}[Modular Fisher sensitivity]\label{def:2.1b}
Fix a self-adjoint element
\(h\in\mathcal{A}_{\mathrm{obs}}\) and define the perturbed family
\[
\omega_s^{(h)}
:=
\omega_{\mathrm{obs}}\circ\mathrm{Ad}(e^{ish}),
\qquad s\in(-\epsilon,\epsilon).
\]
The modular Fisher sensitivity of
\(\omega_{\mathrm{obs}}\) along \(h\) is
\[
\mathcal{I}_F(\omega_{\mathrm{obs}},h)
:=
\mathbf{Q}_{\omega_{\mathrm{obs}}}\!\left(
\frac{d\omega_s^{(h)}}{ds}\bigg|_{s=0},\;
\frac{d\omega_s^{(h)}}{ds}\bigg|_{s=0}
\right).
\]
The modular Fisher sensitivity of the observer is
\[
\mathcal{I}_F(\omega_{\mathrm{obs}})
:=
\sup\bigl\{
\mathcal{I}_F(\omega_{\mathrm{obs}},h)
:\;
h=h^*\in\mathcal{A}_{\mathrm{obs}},\;
\|h\|\le 1
\bigr\}.
\]
\end{definition}

The supremum is taken over unit-norm self-adjoint perturbations in the
observer's own accessible algebra. It is finite whenever the BKM form
is bounded on the unit ball of
\(\mathcal{A}_{\mathrm{obs}}\)-tangent directions, which holds in the
local regime used throughout this paper. The quantity
\(\mathcal{I}_F(\omega_{\mathrm{obs}})\) measures the observer's
maximal intrinsic sensitivity to unitary perturbations of its own
accessible state. It is strictly positive whenever
\(\omega_{\mathrm{obs}}\) is not a trace state on
\(\mathcal{A}_{\mathrm{obs}}\).

Throughout the paper, the term \emph{local modular drift family}
refers to the one-parameter family \(\{\omega_s^{(h)}\}\) for a fixed
perturbation \(h\). When the supremum in
Definition~\ref{def:2.1b} is attained, we take \(h\) to be a
maximizer; when it is not, we use a standard
\(\eta\)-optimization: for each \(\eta>0\), choose
\(h_\eta=h_\eta^*\in\mathcal{A}_{\mathrm{obs}}\) with
\(\|h_\eta\|\le 1\) and
\(\mathcal{I}_F(\omega_{\mathrm{obs}},h_\eta)\ge
\mathcal{I}_F(\omega_{\mathrm{obs}})-\eta\), carry out the proof for
\(h_\eta\), and let \(\eta\downarrow 0\). Despite the name, this
family is not the Tomita--Takesaki modular orbit (which is
state-preserving) but the accessible-unitary perturbation family that
serves as its operational proxy at the level of Fisher information.
This is the family along which the observer's self-description is
hardest to compress.

\subsection{Admissible Self-Models}\label{subsec:2.2}

The theorem chain does not require a literal embedding of the self-model
into a matrix algebra. That stronger route would force us to solve a
much harder operator-algebraic problem than the present paper needs.
What matters here is the existence of a finite statistical bottleneck.

\begin{definition}[Admissible self-model]\label{def:2.2}
An admissible \(n\)-dimensional self-model is a pair
\[
(\Pi_n,\Psi_n),
\]
where:
\begin{enumerate}
\item \(\Pi_n:\mathcal{S}(\mathcal{A}_{\mathrm{obs}})\to
\mathfrak{M}_n\) is the state-space map induced by a normal unital
completely positive map from a finite-dimensional algebra into
\(\mathcal{A}_{\mathrm{obs}}\), so that Petz monotonicity of relative
entropy applies;
\item \(\Psi_n:\mathfrak{M}_n\to
\mathcal{S}(\mathcal{A}_{\mathrm{obs}})\) is a recovery map;
\item \(\Pi_n\) factors through the observer's own accessible state
space, so no external ``God's-eye'' parameter is introduced;
\item near the reference state \(\omega_{\mathrm{obs}}\), the Fisher
matrix on \(\mathfrak{M}_n\) can be whitened:
\[
\mathbf{I}^{(n)}_\omega
=
\mathrm{diag}(\lambda_1,\dots,\lambda_n),
\qquad
\lambda_a \leq \mathcal{I}_F(\omega_{\mathrm{obs}}).
\]
\end{enumerate}
\end{definition}

The last condition is the substantive standing assumption of the
theorem, not a mere coordinate normalization. It says that the
compressed model has only \(n\) effective Fisher directions, and none of
them is sharper than the observer's own modular sensitivity scale. This
is the exact place where the theorem absorbs finite-dimensionality.

\subsection{Self-Description Deficit}\label{subsec:2.3}

The correct loss function is not a metric but a divergence. In a
Type~III setting there is no preferred trace-class density-matrix
picture, while Araki relative entropy is native to the operator
algebraic formulation \cite{Araki1976}.

\begin{definition}[Self-description deficit]\label{def:2.3}
Define the BKM ball of radius \(\delta\) around
\(\omega_{\mathrm{obs}}\) by
\[
\mathcal{B}_\delta(\omega_{\mathrm{obs}})
:=
\bigl\{
\omega\in\mathcal{S}(\mathcal{A}_{\mathrm{obs}})
:\;
\mathbf{Q}_{\omega_{\mathrm{obs}}}(\omega-\omega_{\mathrm{obs}},\,
\omega-\omega_{\mathrm{obs}})
<\delta^2
\bigr\}.
\]
Fix once and for all a \(\delta>0\) small enough that Araki relative
entropy admits its quadratic expansion on
\(\mathcal{B}_\delta(\omega_{\mathrm{obs}})\), and large enough that
the ego-rigidity prior \(\pi_{\mathrm{ego}}\) of
T-DOME~Paper~III\cite{TDOMEIII} can be supported on a segment of the
local modular drift family inside
\(\mathcal{B}_\delta(\omega_{\mathrm{obs}})\). The self-description
deficit of complexity \(n\) is
\[
\varepsilon(n)
:=
\inf_{(\Pi_n,\Psi_n)}
\sup_{\omega\in \mathcal{B}_\delta(\omega_{\mathrm{obs}})}
S\!\left(
\omega
\;\middle\|\;
\Psi_n\circ\Pi_n(\omega)
\right),
\]
where the infimum is taken over all admissible \(n\)-dimensional
self-models.
\end{definition}

Near a fixed state, Araki relative entropy has the familiar BKM/QFIM
quadratic expansion
\[
S(\omega\|\omega+\delta\omega)
=
\frac{1}{2}\mathbf{Q}_\omega(\delta\omega,\delta\omega)
+o(\|\delta\omega\|^2),
\]
which is the local information-geometric structure already used in
Paper~I\cite{PaperI} and in Petz's monotone-metric framework
\cite{Petz1996}. The minimax form is essential: it rules out trivial
constant recovery maps that memorize the reference state while failing
arbitrarily nearby. Thus \(\varepsilon(n)\) measures optimal local
uniform reconstruction loss in the intrinsic Type~III geometry of
states.

\section{The Fisher Ceiling for Admissible Self-Models}\label{sec:3}

\subsection{Petz Monotonicity and Finite Fisher Budget}\label{subsec:3.1}

The first ingredient is standard.

\begin{lemma}[Petz monotonicity]\label{lem:3.1}
For any normal coarse-graining \(\Pi_n\),
\[
S(\omega\|\sigma)\ge S(\Pi_n(\omega)\|\Pi_n(\sigma)).
\]
\end{lemma}

\begin{proof}
This is the monotonicity of relative entropy under coarse-graining
\cite{Petz1986}. The present paper uses it only in the local regime,
where it implies that compression cannot increase BKM/Fisher
information.
\end{proof}

Whitening then converts finite statistical dimension directly into two
distinct ceilings: one for the specific modular drift channel, and one
for the total trace budget across all retained directions.

\begin{proposition}[Directional and total Fisher ceilings]\label{prop:3.2}
Let \((\Pi_n,\Psi_n)\) be an admissible \(n\)-dimensional self-model.
Then:
\begin{enumerate}
\item along the local modular drift family through
\(\omega_{\mathrm{obs}}\), the Fisher information retained after
compression satisfies
\[
\mathcal{I}^{(n)}_{\mathrm{drift}}
\le
\mathcal{I}_F(\omega_{\mathrm{obs}});
\]
\item the total retained data Fisher information satisfies
\[
\mathcal{I}^{(n)}_{\mathrm{data}}
:=
\mathrm{tr}\,\mathbf{I}^{(n)}_\omega
\le
n\,\mathcal{I}_F(\omega_{\mathrm{obs}}).
\]
\end{enumerate}
\end{proposition}

\begin{proof}
\(\!\)(1) Apply Lemma~\ref{lem:3.1} to the one-parameter modular drift
family. In the local quadratic regime, monotonicity of relative entropy
implies monotonicity of the associated BKM/Fisher metric, so
compression cannot increase Fisher information for that specific scalar
parameter.

\(\!\)(2) In Fisher-whitened coordinates, the compressed Fisher matrix is
\[
\mathbf{I}^{(n)}_\omega
=
\mathrm{diag}(\lambda_1,\dots,\lambda_n),
\qquad
\lambda_a \le \mathcal{I}_F(\omega_{\mathrm{obs}}).
\]
Summing over the \(n\) effective directions gives
\[
\mathrm{tr}\,\mathbf{I}^{(n)}_\omega
=
\sum_{a=1}^n \lambda_a
\le
n\,\mathcal{I}_F(\omega_{\mathrm{obs}}).
\]
\end{proof}

\subsection{Ego Rigidity and the Local Deficit Bound}\label{subsec:3.2}

T-DOME~Paper~III\cite{TDOMEIII} introduced a self-referential
calibration loop and the associated Bayesian correction term
\(\mathcal{I}_{\mathrm{ego}}\), interpreted there as ego rigidity. In
the present context the same term enters naturally as prior Fisher
information in a van Trees bound.

\begin{theorem}[Local minimax drift floor]\label{thm:3.3}
Let \((\mathcal{A}_{\mathrm{obs}},\omega_{\mathrm{obs}},\mathfrak{M}_n)\)
be a finite observer, and let \((\Pi_n,\Psi_n)\) be an admissible
\(n\)-dimensional self-model. Assume:
\begin{enumerate}[label=(\roman*)]
\item the local modular drift subfamily
through \(\omega_{\mathrm{obs}}\) lies inside
\(\mathcal{B}_{\delta/2}(\omega_{\mathrm{obs}})\);
\item the quadratic expansion of Araki relative entropy is controlled
on \(\mathcal{B}_\delta(\omega_{\mathrm{obs}})\): there exists
\(\eta\in[0,\tfrac{1}{2})\) such that
\[
S(\omega\|\sigma)
\ge
(\tfrac{1}{2}-\eta)\,
\mathbf{Q}_\omega(\sigma-\omega,\,\sigma-\omega)
\]
for all \(\omega,\sigma\in
\mathcal{B}_\delta(\omega_{\mathrm{obs}})\);
\item the BKM form is uniformly comparable on
\(\mathcal{B}_\delta(\omega_{\mathrm{obs}})\): there exists
\(C_\delta\ge 1\) such that
\[
C_\delta^{-1}\,
\mathbf{Q}_{\omega_{\mathrm{obs}}}(v,v)
\le
\mathbf{Q}_\omega(v,v)
\le
C_\delta\,
\mathbf{Q}_{\omega_{\mathrm{obs}}}(v,v)
\]
for all \(\omega\in
\mathcal{B}_\delta(\omega_{\mathrm{obs}})\) and all tangent vectors
\(v\), with \(C_\delta\to 1\) as \(\delta\to 0\).
\end{enumerate}
Define the \emph{local self-description deficit} by restricting the
minimax to self-models whose reconstruction maps
\(\mathcal{B}_{\delta/2}(\omega_{\mathrm{obs}})\) into
\(\mathcal{B}_\delta(\omega_{\mathrm{obs}})\):
\[
\varepsilon_{\mathrm{loc}}(n)
:=
\inf_{\substack{(\Pi_n,\Psi_n)\\[2pt]
\Psi_n\circ\Pi_n(\mathcal{B}_{\delta/2})\subset
\mathcal{B}_\delta}}
\sup_{\omega\in \mathcal{B}_\delta(\omega_{\mathrm{obs}})}
S\!\left(
\omega
\;\middle\|\;
\Psi_n\circ\Pi_n(\omega)
\right).
\]
Then
\[
\varepsilon_{\mathrm{loc}}(n)
\ge
\frac{(1-2\eta)\,\mathcal{I}_F(\omega_{\mathrm{obs}})}
{2\bigl(\mathcal{I}_F(\omega_{\mathrm{obs}})
+
\mathcal{I}_{\mathrm{ego}}\bigr)}.
\]
In particular, \(\varepsilon(n)=\varepsilon_{\mathrm{loc}}(n)\)
whenever the infimum in Definition~\ref{def:2.3} is attained by a
locally bounded self-model.
\end{theorem}

\begin{proof}
Fix a locally bounded admissible self-model \((\Pi_n,\Psi_n)\), and
write \(\{\omega_s\}\) for the local modular drift family. Let
\(\pi\) be the ego-rigidity prior \(\pi_{\mathrm{ego}}\) restricted
and renormalized to the drift segment inside
\(\mathcal{B}_{\delta/2}(\omega_{\mathrm{obs}})\) (this is possible by
assumption~(i) and the support condition in
Definition~\ref{def:2.3}). Then
\[
\sup_{\omega\in \mathcal{B}_\delta(\omega_{\mathrm{obs}})}
S\!\left(
\omega\;\middle\|\;\Psi_n\circ\Pi_n(\omega)
\right)
\ge
\int
S\!\left(
\omega_s\;\middle\|\;\Psi_n\circ\Pi_n(\omega_s)
\right)\,d\pi(s).
\]
Since \(\omega_s\in\mathcal{B}_{\delta/2}(\omega_{\mathrm{obs}})\)
and the self-model is locally bounded,
\(\Psi_n\circ\Pi_n(\omega_s)\in
\mathcal{B}_\delta(\omega_{\mathrm{obs}})\). Both states therefore
lie inside the controlled ball, so assumption~(ii) gives
\[
S(\omega_s\|\Psi_n\circ\Pi_n(\omega_s))
\ge
(\tfrac{1}{2}-\eta)\,
\mathbf{Q}_{\omega_s}(\delta\omega_{\mathrm{recon}}(s),\,
\delta\omega_{\mathrm{recon}}(s)),
\]
where
\(\delta\omega_{\mathrm{recon}}(s):=
\Psi_n\circ\Pi_n(\omega_s)-\omega_s\).
Since \(\omega_s\in\mathcal{B}_{\delta/2}(\omega_{\mathrm{obs}})\),
assumption~(iii) gives
\(\mathbf{Q}_{\omega_s}(v,v)\ge C_\delta^{-1}\,
\mathbf{Q}_{\omega_{\mathrm{obs}}}(v,v)\). Absorbing
\(C_\delta^{-1}\) into the prefactor \((\tfrac{1}{2}-\eta)\) (both
approach their ideal values as \(\delta\to 0\)), we may replace
\(\mathbf{Q}_{\omega_s}\) by
\(\mathbf{Q}_{\omega_{\mathrm{obs}}}\) at the cost of adjusting
\(\eta\). Decompose the reconstruction error:
\[
\delta\omega_{\mathrm{recon}}(s)
=
\dot{\omega}\,(\widehat{s}(s)-s)
+
\delta\omega_{\perp}(s),
\]
where \(\widehat{s}(s)\) is the
\(\mathbf{Q}_{\omega_{\mathrm{obs}}}\)-projection coefficient of
\(\delta\omega_{\mathrm{recon}}(s)\) onto the drift tangent
\(\dot{\omega}\).
By positivity and the orthogonal decomposition,
\[
\mathbf{Q}_{\omega_{\mathrm{obs}}}(\delta\omega_{\mathrm{recon}}(s),\,
\delta\omega_{\mathrm{recon}}(s))
\ge
\mathcal{I}_F(\omega_{\mathrm{obs}})
\bigl(\widehat{s}(s)-s\bigr)^2.
\]
Integrating against \(\pi\) gives
\[
\int
S\!\left(
\omega_s\;\middle\|\;\Psi_n\circ\Pi_n(\omega_s)
\right)\,d\pi(s)
\ge
(\tfrac{1}{2}-\eta)\,\mathcal{I}_F(\omega_{\mathrm{obs}})
\int\bigl(\widehat{s}(s)-s\bigr)^2\,d\pi(s).
\]

By Proposition~\ref{prop:3.2}(1), the compressed model retains at most
\(\mathcal{I}_F(\omega_{\mathrm{obs}})\) Fisher information about that
single scalar parameter. Adding the Bayesian prior term from ego
rigidity gives the effective scalar information ceiling
\[
\mathcal{I}_{\mathrm{eff}}^{\mathrm{drift}}
\le
\mathcal{I}_F(\omega_{\mathrm{obs}})
+
\mathcal{I}_{\mathrm{ego}}.
\]
The scalar van Trees inequality then yields a Bayes-risk lower bound
\[
\int\bigl(\widehat{s}(s)-s\bigr)^2\,d\pi(s)
\ge
\frac{1}{\mathcal{I}_F(\omega_{\mathrm{obs}})
+
\mathcal{I}_{\mathrm{ego}}}
\]
for estimation along the local modular drift channel
\cite{vanTrees1968,TDOMEIII}. Combining gives
\[
\sup_{\omega\in \mathcal{B}_\delta(\omega_{\mathrm{obs}})}
S\!\left(
\omega\;\middle\|\;\Psi_n\circ\Pi_n(\omega)
\right)
\ge
\frac{(1-2\eta)\,\mathcal{I}_F(\omega_{\mathrm{obs}})}
{2\bigl(\mathcal{I}_F(\omega_{\mathrm{obs}})
+
\mathcal{I}_{\mathrm{ego}}\bigr)}.
\]
Since this holds for every locally bounded admissible
\((\Pi_n,\Psi_n)\), the same lower bound holds after taking the
infimum.
\end{proof}

\begin{corollary}[Non-vanishing local floor]\label{cor:3.4}
If \(\mathcal{I}_F(\omega_{\mathrm{obs}})>0\) and \(\eta\) can be
taken arbitrarily small (by shrinking \(\delta\)), then for every
\(n\),
\[
\varepsilon_{\mathrm{loc}}(n)\ge
\frac{1}{
2\left(
1+\mathcal{I}_{\mathrm{ego}}/\mathcal{I}_F(\omega_{\mathrm{obs}})
\right)}.
\]
\end{corollary}

\begin{proof}
Take \(\eta\to 0\) in Theorem~\ref{thm:3.3} and rewrite
algebraically.
\end{proof}

In words: the local self-description deficit is bounded below by a
ratio that depends only on the observer's modular sensitivity and its
ego-rigidity prior. Even if \(\mathcal{I}_{\mathrm{ego}}=0\), the
floor is \(1/2\), reflecting a pure Fisher ceiling. A larger
\(\mathcal{I}_{\mathrm{ego}}\) lowers the floor because the prior
absorbs some of the reconstruction burden, but the floor never reaches
zero as long as \(\mathcal{I}_F>0\).

\begin{remark}\label{rem:3.5}
Theorem~\ref{thm:3.3} is the theorem-grade core of Paper~III, but its
scope is deliberately narrow. It is a local minimax statement for the
single modular drift channel around a fixed observer state, restricted
to locally bounded self-models. The
\(n\)-dependence of higher-order control enters only later, through the
full trace budget of Proposition~\ref{prop:3.2}(2) applied to
\(\mathbf{Q}_\omega\)-orthogonal coefficient families. A stronger global
operator-algebraic route is left to future work.
\end{remark}

\section{Small-Ball Remainders and Finite Control Order}\label{sec:4}

\subsection{From Paper II to a Control-Order Definition}\label{subsec:4.1}

Paper~II\cite{PaperII} showed that, in the CHM small-ball framework
\cite{CasiniHuertaMyers2011},
\[
\delta\langle K_{B,u}\rangle
=
\sum_{k=0}^{\infty} a_k(x_0,u)\,R^{d+1+2k},
\]
with the \(k=0\) term fixed by the local stress tensor and the first
uncontrolled terms appearing at \(O(R^{d+3})\). To turn this into a
finite-self-model statement, we now leave the single-channel setting of
Theorem~\ref{thm:3.3} and use the full trace budget from
Proposition~\ref{prop:3.2}(2). We associate a Fisher cost to each
higher-order coefficient.

\begin{definition}[Coefficient Fisher information]\label{def:4.1}
Fix \(k\ge 1\). Let \(\omega_\theta^{(k)}\) be a one-parameter family of
perturbations whose tangent vector is chosen, by BKM Gram--Schmidt
orthogonalization if necessary, to be
\(\mathbf{Q}_\omega\)-orthogonal to the span of the lower-order
coefficient directions and to vary \(a_k\) at linear order. The Fisher
information of the \(k\)-th coefficient is
\[
I_k
:=
\mathbf{Q}_\omega\!\left(
\frac{\partial\omega_\theta^{(k)}}{\partial\theta}\Big|_{\theta=0},
\frac{\partial\omega_\theta^{(k)}}{\partial\theta}\Big|_{\theta=0}
\right).
\]
\end{definition}

We then encode finite control by budget exhaustion.

\begin{definition}[Small-ball control order]\label{def:4.2}
The maximal small-ball control order of an \(n\)-dimensional self-model
is
\[
k(n)
:=
\max\left\{
K\ge 1:
\sum_{j=1}^{K} I_j^{-1}
\le
n\,\mathcal{I}_F(\omega_{\mathrm{obs}})
+
\mathcal{I}_{\mathrm{ego}}
\right\}.
\]
\end{definition}

This definition packages the additive Fisher-budget approximation
directly into the quantity being studied. The point is not that all
channels are literally independent in every model, but that the joint
control problem costs at least the cumulative Cram\'er--Rao budget
appearing on the left-hand side. With the orthogonalized coefficient
directions of Definition~\ref{def:4.1}, the local Fisher matrix is
diagonal at the reference state, so the additive budget written above is
exact at quadratic order.

\subsection{Abstract Growth Law}\label{subsec:4.2}

The next statement is conditional on the spectral asymptotic extracted
from the worked example in Section~\ref{sec:5}.

\begin{theorem}[Finite self-models control only finitely many orders
(conditional on spectral hypothesis)]
\label{thm:4.3}
Assume the higher-order coefficient spectrum obeys
\[
\log I_k = -\alpha k\log k + O(k)
\qquad
(\alpha>0).
\]
Then
\[
k(n)
=
O\!\left(\frac{\log n}{\log\log n}\right).
\]
\end{theorem}

\begin{proof}
By Definition~\ref{def:4.2}, which packages the trace ceiling from
Proposition~\ref{prop:3.2}(2),
\[
\sum_{j=1}^{k(n)} I_j^{-1}
\le
n\,\mathcal{I}_F(\omega_{\mathrm{obs}})
+
\mathcal{I}_{\mathrm{ego}}.
\]
The assumed asymptotic implies that \(I_j^{-1}\) grows in a
stretched-factorial manner, so the sum is dominated by its final term:
\[
\sum_{j=1}^{K} I_j^{-1}\asymp I_K^{-1}.
\]
Taking logarithms gives
\[
\alpha K\log K + O(K)
\lesssim
\log n.
\]
Solving this growth relation yields
\[
K = O\!\left(\frac{\log n}{\log\log n}\right).
\]
\end{proof}

\begin{remark}\label{rem:4.4}
Theorem~\ref{thm:4.3} says less than a general theorem on all higher
coefficients, but more than a slogan. It shows that once the spectrum is
super-polynomially sparse, increasing model complexity produces only
logarithmic gains in controllable small-ball order.
\end{remark}

\section{Worked Example: Free Scalar CFT and the Fisher Spectrum}
\label{sec:5}

\subsection{What Must Be Verified}\label{subsec:5.1}

The only input missing from Section~\ref{sec:4} is the growth law for
\(I_k\). Here we do not attempt a general theorem for all conformal
theories. We verify a worked-example-level statement in the free scalar
CFT, which is enough to support the first version of Paper~III.

The logic is simple. The Fisher spectrum of the \(k\)-th coefficient is
controlled by two pieces:
\begin{enumerate}
\item a signal factor \(c_k\) coming from the CHM moment extracting the
\(2k\)-th derivative;
\item a renormalized boundary kernel \(\gamma_k\) coming from the
fluctuation structure of the smeared operator.
\end{enumerate}
If \(\gamma_k\) contributes at most polynomial \(k\)-dependence and no
new factorial growth, then the asymptotic is governed entirely by the
signal factor.

\subsection{Exact d=1 Signal Factor}\label{subsec:5.2}

In \(d=1\), the CHM modular Hamiltonian is
\[
K_B
=
\frac{\pi}{R}\int_{-R}^{R}(R^2-x^2)\,T_{00}(x)\,dx.
\]
Expanding \(T_{00}(x)\) around the ball center gives
\[
T_{00}(x)
=
\sum_{k\ge 0}\frac{x^{2k}}{(2k)!}\partial_x^{2k}T_{00}(0)+\cdots.
\]
Therefore the coefficient multiplying the \(2k\)-th derivative is
\[
c_k^{(d=1)}
=
\frac{\pi}{R}\cdot\frac{1}{(2k)!}
\int_{-R}^{R}(R^2-x^2)x^{2k}\,dx\cdot R^{-(2k+2)}.
\]
The radial moment is exact:
\[
\int_{-R}^{R}(R^2-x^2)x^{2k}\,dx
=
\frac{4R^{2k+3}}{(2k+1)(2k+3)}.
\]
Hence
\[
c_k^{(d=1)}
=
\frac{4\pi}{(2k+1)(2k+3)(2k)!},
\]
which shows exact factorial suppression:
\[
c_k^{(d=1)}
\sim
\frac{\mathrm{const}}{(2k)!\,k^2}.
\]

\subsection{Boundary Kernel}\label{subsec:5.3}

The kernel side is less rigid, and this is exactly where we refuse to
overclaim. The large-\(k\) boundary-layer analysis shows that in
\(d=1\),
\[
\gamma_k^{(d=1)}\sim \frac{G_*}{4k^4},
\]
so no factorial enhancement appears. In higher dimensions, the angular
average is ultraviolet singular before renormalization, and the correct
object is a renormalized boundary kernel. What the present paper needs,
and what the worked example supports, is only the following structural
input:

\begin{proposition}[Worked-example spectral input]\label{prop:5.1}
In the free scalar CFT worked example:
\begin{enumerate}
\item the \(d=1\) signal factor is exactly factorially suppressed, as
above;
\item if the renormalized boundary kernel contributes at most
polynomial \(k\)-dependence and no new factorial growth in general
dimension, then the coefficient Fisher spectrum satisfies
\[
\log I_k = -\alpha k\log k + O(k)
\qquad
(\alpha>0).
\]
\end{enumerate}
\end{proposition}

\begin{proof}
The first part is the exact computation above. For the second, the
polynomial kernel contributes only \(O(\log k)\) corrections to
\(\log I_k\), while the signal factor contributes the dominant
\(-\alpha k\log k\) term through Stirling asymptotics of \((2k)!\).
\end{proof}

\begin{corollary}\label{cor:5.2}
Under the worked-example hypothesis of
Proposition~\ref{prop:5.1}, the small-ball control order satisfies
\[
k(n)=O\!\left(\frac{\log n}{\log\log n}\right).
\]
\end{corollary}

\begin{proof}
Combine Proposition~\ref{prop:5.1} with
Theorem~\ref{thm:4.3}.
\end{proof}

\subsection{What the Worked Example Shows and Does Not Show}
\label{subsec:5.4}

The worked example shows three things. First, the spectral input needed
for Theorem~\ref{thm:4.3} is physically plausible and mathematically
supported in the simplest CFT setting. Second, the dominant suppression
comes from the CHM signal factor rather than from a delicate estimate of
the boundary kernel exponent. Third, the resulting growth law is slower
than the original power-law guess: finite self-models gain access to
higher-order information only logarithmically.

What it does \emph{not} show is equally important. It does not prove a
general-dimensional closed formula for \(\gamma_k\). It does not prove
a universal theorem for all conformal nets. And it does not yet convert
the spectral statement into a purely operator-algebraic result
independent of the worked-example assumptions.

\section{Trilogy Implications}\label{sec:6}

\subsection{Quantifying Paper G}\label{subsec:6.1}

Paper~G\cite{PaperG} argued that the explanatory boundary is structural
rather than epistemic. Theorem~\ref{thm:3.3} gives that claim a first
quantitative form. Even before one studies higher-order coefficient
families, every locally bounded admissible finite self-model with
\(\mathcal{I}_F(\omega_{\mathrm{obs}})>0\) has a strictly positive local
minimax loss along the observer's own modular drift channel. The limit
is not merely that an observer happens to lack some information; it is
that finite self-description already hits a nonzero local floor,
provided the self-model's reconstruction remains in the quadratic
regime.

\subsection{Relocating the Gap in Paper II Terms}\label{subsec:6.2}

Paper~II\cite{PaperII} showed that linearized gravitational consistency
fixes the \(R^{d+1}\) term in the small-ball expansion, leaving the
higher-order remainder outside the linearized control problem. Paper~III
now reinterprets that remainder. The uncontrolled orders are not just
``higher corrections'' in a formal expansion. They are the orders whose
joint Fisher cost exceeds what a finite self-model can stably absorb.

This is the precise sense in which the gap is relocated. It is not a
gap between two ontologically separate domains. It is the mismatch
between a finite internal model and the higher-order structure of its
own modular-geometric organization.

\subsection{Trilogy-Level Interpretation}\label{subsec:6.3}

Conditional on the identifications defended in Papers~I and~II
\cite{PaperI,PaperII}, the hard-problem headline can now be restated
more carefully. Paper~I removed the presupposition of independent
describability. Paper~II showed that, in higher dimensions, consistency
of the inseparable structures forces linearized gravity. Paper~III adds
that the residual sense of an explanatory gap is quantitatively
compatible with a finite self-description horizon.

We do not identify qualia with specific coefficient orders or with a
particular small-ball remainder tower. The present claim is narrower:
once the trilogy's modular and geometric identifications are granted,
the remaining gap is structurally and quantitatively located in the part
of the algebraic organization that finite self-models cannot stably
internalize.

\section{Limitations and Open Problems}\label{sec:7}

The current paper leaves four substantial tasks open.

\begin{enumerate}
\item \textbf{Strong Lemma C.} The Fisher ceiling is currently proved
for self-models that factor through a Fisher-whitened finite statistical
model. A stronger theorem would derive the same ceiling directly from
pure operator-algebraic self-reference, without inserting that
statistical layer by hand. The obstacle is that the passage from a Type
III von Neumann algebra to a finite-dimensional statistical model
currently requires the admissibility assumption
(Definition~\ref{def:2.2}), which is not derived from the algebraic
structure.

\item \textbf{General-dimensional kernel control.} The free scalar
worked example supports the spectral law under a ``no new factorial
growth'' assumption on the renormalized boundary kernel. Turning that
assumption into a theorem remains a real mathematical problem. The
obstacle is that the renormalized boundary kernel in \(d>1\) involves
UV-singular angular integrals whose large-\(k\) behavior is not
controlled by the present methods.

\item \textbf{Global and nonlinear extension.} Paper~II was deliberately
linearized and local. A true sequel would connect the present finite
self-description bound to global ball-space consistency and to the
nonlinear regime.

\item \textbf{Observer design and optimal rigidity.} T-DOME~Paper~III
already suggests a rigidity-sensitivity trade-off in
\(\mathcal{I}_{\mathrm{ego}}\). Paper~III does not optimize that trade
off; it only shows how it enters the deficit bound. The problem of
optimal self-description at fixed thermodynamic and computational budget
is still open.
\end{enumerate}

\section{Conclusion}\label{sec:8}

Paper~III provides a first quantitative pass over the explanatory
gap. Paper~G located the gap architecturally. Paper~II located a
physical carrier of the unresolved information in higher-order
small-ball terms. The present paper introduces a finite-self-model
formalism in which that residual information becomes a measurable
deficit.

Within the deliberately narrow regime studied here, the conclusion is
sharp. A locally bounded finite self-model has a finite Fisher ceiling.
Along the observer's own modular drift channel, that ceiling produces a
nonzero local minimax loss whenever
\(\mathcal{I}_F(\omega_{\mathrm{obs}})>0\), provided the
reconstruction remains in the quadratic regime. Across
\(\mathbf{Q}_\omega\)-orthogonal higher-order small-ball directions, it
allows only logarithmically slow access to new orders. The gap is
therefore not eliminated, but neither is it mysterious. It appears as a
structural horizon imposed by finite self-description.

\appendix

\section{Trilogy Key Premises Summary}\label{app:premises}

The following table lists the specific results from other papers in the
trilogy that are used in Paper~III. Each entry provides a self-contained
statement sufficient for following the proofs in the present paper.

\begin{enumerate}
\item \textbf{BKM/QFIM structure} (Paper~I\cite{PaperI}).
The Araki relative entropy of two nearby faithful normal states on a
Type~III von Neumann algebra admits a local quadratic expansion
\(S(\omega\|\omega+\delta\omega)=\tfrac{1}{2}\mathbf{Q}_\omega
(\delta\omega,\delta\omega)+o(\|\delta\omega\|^2)\), where
\(\mathbf{Q}_\omega\) is the BKM/QFIM quadratic form. This is the
information-geometric structure used in Definitions~\ref{def:2.1b}
and~\ref{def:2.3} and throughout Section~\ref{sec:3}.

\item \textbf{Small-ball expansion with \(O(R^{d+3})\) remainder}
(Paper~II\cite{PaperII}).
In the CHM ball framework for conformal nets, the modular energy of a
small ball of radius \(R\) admits the expansion
\(\delta\langle K_{B,u}\rangle = \frac{2\pi V_d}{d+2}R^{d+1}
\delta\langle T_{\mu\nu}\rangle u^\mu u^\nu + O(R^{d+3})\).
The \(k=0\) term is fixed by the linearized Einstein equation; the
higher-order terms are uncontrolled. This expansion is the starting
point of Section~\ref{sec:4}.

\item \textbf{Qualitative incompleteness} (Paper~G\cite{PaperG}).
No accessibility-based framework can fully internalize its own
accessibility conditions without either regress or collapse: the
self-referential fixed point
\(\mathcal{A}^*\in\mathrm{Acc}(\mathcal{A}^*)\) generically does not
exist. This is the qualitative claim that Paper~III quantifies.

\item \textbf{Ego-rigidity prior}
(T-DOME~Paper~III\cite{TDOMEIII}).
The ego-rigidity prior \(\mathcal{I}_{\mathrm{ego}}\) is the Fisher
information of the observer's prior distribution over its own
self-model parameters:
\(\mathcal{I}_{\mathrm{ego}}:=\int(\partial_\sigma\log
\pi_{\mathrm{ego}}(\sigma))^2\pi_{\mathrm{ego}}(\sigma)\,d\sigma\).
It enters Theorem~\ref{thm:3.3} as the Bayesian prior term in the van
Trees bound.
\end{enumerate}
 
\newpage
%
\addcontentsline{toc}{section}{References}

\begin{thebibliography}{99}

\bibitem{PaperI}
S.~Liu,
\emph{The Inseparability of Geometry and Modular Flow: A Structural Inseparability Theorem and Its Implications for the Hard Problem},
Consciousness Trilogy, Paper~I (2026),
Zenodo, DOI: 10.5281/zenodo.19607387.

\bibitem{PaperII}
S.~Liu,
\emph{Higher-Dimensional Consistency of Geometry and Modular Flow: Small-Ball Modular Data and Linearized Gravity},
Consciousness Trilogy, Paper~II (2026),
Zenodo, DOI: 10.5281/zenodo.19607543.

\bibitem{PaperG}
S.~Liu,
\emph{Structural Limits of Unification: Accessibility, Incompleteness, and the Necessity of a Final Cut},
HAFF, Paper~G (2026),
Zenodo, DOI: 10.5281/zenodo.18402907.

\bibitem{TDOMEIII}
S.~Liu,
\emph{Thermodynamic Dynamics of Observer-Memory Entanglement, Paper~III: The Self-Calibration Loop},
T-DOME, Paper~III (2026),
Zenodo, DOI: 10.5281/zenodo.18591771.

\bibitem{Araki1976}
H.~Araki,
\emph{Relative entropy of states of von Neumann algebras},
Publ.\ Res.\ Inst.\ Math.\ Sci.\ \textbf{11}, 809--833 (1976).

\bibitem{Petz1986}
D.~Petz,
\emph{Sufficient subalgebras and the relative entropy of states of a von Neumann algebra},
Commun.\ Math.\ Phys.\ \textbf{105}, 123--131 (1986).

\bibitem{Petz1996}
D.~Petz,
\emph{Monotone metrics on matrix spaces},
Linear Algebra Appl.\ \textbf{244}, 81--96 (1996).

\bibitem{vanTrees1968}
H.~L.~Van Trees,
\emph{Detection, Estimation, and Modulation Theory, Part~I},
Wiley, New York (1968).

\bibitem{CasiniHuertaMyers2011}
H.~Casini, M.~Huerta, and R.~C.~Myers,
\emph{Towards a derivation of holographic entanglement entropy},
JHEP \textbf{1105}, 036 (2011),
arXiv:1102.0440.

\end{thebibliography}
 
\end{document}
